% Requires \usepackage{amsmath}. Place before the "Species distribution models"
% subsection; then the short connectivity paragraph there can be trimmed to avoid
% duplication. Replace citation keys with your .bib entries.

\subsection{Larval connectivity metrics}

Larval dispersal was simulated with a biophysical individual-based model that
couples a hydrodynamic ocean model with Lagrangian particle tracking and a
larval pelagic-duration/behaviour module \citep{yourABMref}. The study area was
discretized into $n=2340$ grid cells. Virtual larvae were released from each
cell and tracked until settlement; aggregating over all release events yielded a
connectivity matrix $C$, where $C_{jk}$ is the number of larvae released in cell
$j$ that settled in cell $k$, together with a release vector $R_j$ giving the
number of larvae released from cell $j$.

We converted these into a flow (transition) matrix $F$ of settlement
probabilities,
\begin{equation}
F_{jk} = \frac{C_{jk}}{R_j},
\qquad 0 \le \sum_{k} F_{jk} \le 1,
\end{equation}
so that $F_{jk}$ is the probability that a larva released in source cell $j$
settles in sink cell $k$ (rows index sources, columns sinks). We treated $F$ as
the weighted adjacency matrix of a directed dispersal network whose nodes are
grid cells and whose edge weights are settlement probabilities.

From this network we derived cell-level connectivity metrics. Because they depend
only on the dispersal network and not on the observed cockle densities, these
metrics are non-circular predictors of cockle distribution. For each cell $k$ we
computed the weighted in-strength, the local retention, and the eigenvector
centrality,
\begin{align}
I_k &= \sum_{j} F_{jk}, \\
L_k &= F_{kk}, \\
v_k &: \; \mathbf{v}\ \text{the leading eigenvector of}\ F,\ \ \max_k v_k = 1,
\end{align}
where the in-strength $I_k$ is the total incoming settlement probability (larval
import, or sink strength), the local retention $L_k$ is the probability of
self-recruitment, and the eigenvector centrality $v_k$ ranks cells by their
position in the dispersal network, assigning high values to cells that receive
larvae from other well-connected cells.

In addition to these structural (potential-connectivity) metrics, we computed a
\emph{realized} larval supply that weights each incoming connection by the
predicted biomass of its source cell,
\begin{equation}
S_k = \sum_{j} \hat{b}_j\, F_{jk},
\end{equation}
where $\hat{b}_j$ is the mean biomass of source cell $j$ predicted from a
depth-only spatial model (Section~XX). Because $\hat{b}_j$ is derived from the
cockle data, $S_k$ quantifies realized rather than potential supply; we used the
transform $\log(1+S_k)$ to compress biomass hot-spots while keeping empty source
cells at zero.

Each survey location was assigned the connectivity metrics of the grid cell
containing it. In the species distribution models
(Eq.~\ref{eq:linpred}) the connectivity covariate $c_i$ is the in-strength
$I_{k(i)}$ of the cell $k(i)$ containing observation $i$; we used the structural
in-strength as the primary connectivity metric and compared it against
eigenvector centrality and log realized supply $\log(1+S_k)$ as a sensitivity
analysis. Structural graph metrics were computed with the \texttt{igraph}
package \citep{csardi2006}.
